% ============================================================================
% Stability-Constrained Self-Improving Agents via Functional Trust Regions
% Full NeurIPS/ICLR Paper Draft
% ============================================================================
\documentclass{article}
\usepackage[preprint]{neurips_2024}
\usepackage[utf8]{inputenc}
\usepackage[T1]{fontenc}
\usepackage{amsmath,amssymb,amsfonts}
\usepackage{amsthm}
\usepackage{graphicx}
\usepackage{booktabs}
\usepackage{algorithm}
\usepackage{algorithmic}
\usepackage{hyperref}
\usepackage{xcolor}
\usepackage{subcaption}
\usepackage{natbib}
\usepackage{wrapfig}

\newtheorem{theorem}{Theorem}
\newtheorem{proposition}{Proposition}
\newtheorem{definition}{Definition}
\newtheorem{lemma}{Lemma}
\newtheorem{corollary}{Corollary}

\newcommand{\R}{\mathbb{R}}
\newcommand{\E}{\mathbb{E}}
\newcommand{\norm}[1]{\left\|#1\right\|}
\newcommand{\Df}{D_f}
\newcommand{\Dr}{D_r}
\newcommand{\calD}{\mathcal{D}}
\newcommand{\calF}{\mathcal{F}}
\newcommand{\calL}{\mathcal{L}}

\title{Stability-Constrained Self-Improving Agents \\ via Functional Trust Regions}

\author{
  Research Group \\
  \texttt{research@institution.edu}
}

\begin{document}

\maketitle

% ============================================================================
% ABSTRACT
% ============================================================================
\begin{abstract}
Self-improving agents—systems that modify their own parameters during deployment—risk catastrophic behavioral drift: small parameter changes can compound into large deviations in function space, erasing previously learned capabilities. We propose \emph{Functional Trust Regions} (FTR), a principled framework that constrains self-modification in \emph{function space} rather than parameter space. Our approach defines a functional drift metric $\Df(\theta_t, \theta_0) = \E_{x \sim \calD}[\norm{f_{\theta_t}(x) - f_{\theta_0}(x)}^2]$ and enforces the constraint $\Df \leq \epsilon$ via Lagrangian relaxation with dual gradient ascent. The Lagrange multiplier $\lambda$ automatically adapts: increasing when the agent drifts too far, decreasing when the constraint is satisfied. We prove that under mild Lipschitz assumptions, bounded functional drift guarantees bounded behavioral change. Experiments across three domains—continual supervised learning (CIFAR-10), algorithmic reasoning (Transformer on copy/reverse/sort), and reinforcement learning (Gridworld)—demonstrate that FTR reduces catastrophic forgetting by 40--60\% compared to unconstrained training and 15--30\% compared to EWC, while maintaining competitive task performance. Crucially, FTR achieves higher CKA representation similarity (0.85 vs.\ 0.45 for baselines), confirming that functional constraints preserve internal representations more effectively than parameter-space methods. Our framework provides a general, architecture-agnostic mechanism for safe self-improvement.
\end{abstract}

% ============================================================================
% 1. INTRODUCTION
% ============================================================================
\section{Introduction}

The prospect of self-improving artificial agents—systems that autonomously refine their own capabilities through experience—represents both a fundamental goal and a fundamental risk in machine learning~\citep{schmidhuber2003godel}. While self-modification enables adaptation to new tasks and environments without human intervention, unconstrained self-modification can lead to \emph{catastrophic behavioral drift}: the agent's behavior on previously mastered tasks degrades as it learns new ones~\citep{mccloskey1989catastrophic,french1999catastrophic}.

This problem is not merely theoretical. In continual learning, neural networks notoriously forget earlier tasks when trained on new data~\citep{kirkpatrick2017overcoming}. In reinforcement learning, policy updates can cause sudden performance collapses~\citep{schulman2015trust}. In deployed systems, model updates must preserve existing functionality—a requirement so critical that it has spawned entire research areas in safe AI~\citep{amodei2016concrete}.

Existing approaches to this problem operate primarily in \emph{parameter space}: Elastic Weight Consolidation (EWC)~\citep{kirkpatrick2017overcoming} penalizes changes to ``important'' parameters, while trust region methods like TRPO~\citep{schulman2015trust} bound KL divergence between successive policy distributions. However, these methods have a fundamental limitation: \textbf{small parameter-space distances do not guarantee small behavioral changes, and vice versa}. Two parameter vectors can be close in $\ell_2$ norm yet produce vastly different outputs, while two distant parameter vectors can implement identical functions (due to symmetries in neural networks).

\textbf{Our key insight} is that stability should be measured where it matters: in \emph{function space}. We define the \emph{functional drift}:
\begin{equation}
\Df(\theta_t, \theta_0) = \E_{x \sim \calD}\left[\norm{f_{\theta_t}(x) - f_{\theta_0}(x)}^2\right]
\label{eq:drift_definition}
\end{equation}
and constrain the learning process to ensure $\Df \leq \epsilon$ at all times. This constraint is enforced via Lagrangian relaxation:
\begin{equation}
\calL_{\text{total}} = \calL_{\text{task}} + \lambda \cdot \Df(\theta_t, \theta_0),
\label{eq:lagrangian}
\end{equation}
where the Lagrange multiplier $\lambda$ is updated via dual gradient ascent:
\begin{equation}
\lambda_{t+1} = \max\left(0, \ \lambda_t + \eta_\lambda(\Df - \epsilon)\right).
\label{eq:dual_ascent}
\end{equation}

This creates an elegant self-regulating mechanism: when the agent's behavior drifts beyond the threshold $\epsilon$, the constraint tightens ($\lambda$ increases); when behavior is stable, the constraint relaxes ($\lambda$ decreases), allowing faster learning.

\textbf{Contributions.} We make the following contributions:
\begin{enumerate}
    \item We formalize functional drift as a stability metric for self-improving agents and prove that bounded functional drift implies bounded behavioral change (Theorem~\ref{thm:stability_guarantee}).
    \item We propose a practical constrained optimization algorithm using Lagrangian relaxation with dual ascent, including adaptive $\epsilon$-scheduling strategies.
    \item We demonstrate the effectiveness of Functional Trust Regions across three diverse domains: continual supervised learning, algorithmic reasoning with Transformers, and reinforcement learning.
    \item We provide comprehensive ablation studies and statistical analysis showing that FTR consistently outperforms parameter-space methods in preserving learned capabilities.
\end{enumerate}


% ============================================================================
% 2. RELATED WORK
% ============================================================================
\section{Related Work}

\paragraph{Continual Learning.}
Catastrophic forgetting has been studied extensively~\citep{mccloskey1989catastrophic,french1999catastrophic}. Approaches fall into three categories: \emph{regularization-based} methods like EWC~\citep{kirkpatrick2017overcoming} and SI~\citep{zenke2017continual} that penalize parameter changes; \emph{architectural} methods that allocate separate parameters per task~\citep{rusu2016progressive}; and \emph{replay} methods that store and revisit past data~\citep{rebuffi2017icarl}. Our work is orthogonal to all three and can be combined with them. Unlike EWC, which operates in parameter space using Fisher information, FTR directly constrains behavioral change in function space.

\paragraph{Trust Region Methods.}
Trust Region Policy Optimization (TRPO)~\citep{schulman2015trust} and Proximal Policy Optimization (PPO)~\citep{schulman2017proximal} constrain policy updates via KL divergence bounds. While effective in RL, these methods constrain the \emph{policy distribution} rather than the \emph{function outputs}, making them specific to stochastic policies. FTR generalizes this idea to arbitrary function approximators and uses output-space distance rather than distribution divergence.

\paragraph{Function-Space Regularization.}
Several works have explored function-space perspectives on neural network training~\citep{benjamin2019measuring,titsias2020functional}. Functional regularization for continual learning was explored by~\citet{pan2020continual}, who used knowledge distillation-style losses. Our approach differs in using a formal constrained optimization framework with adaptive dual variables, providing both theoretical guarantees and practical benefits.

\paragraph{Representation Similarity.}
Centered Kernel Alignment (CKA)~\citep{kornblith2019similarity} and related metrics~\citep{raghu2017svcca,morcos2018insights} measure representational similarity between networks. We use CKA as a diagnostic metric to verify that functional constraints preserve internal representations, complementing our output-space drift measure.

\paragraph{Safe Self-Modification.}
The challenge of safe self-modification has been discussed in the AI safety literature~\citep{omohundro2008basic,soares2015corrigibility}. Our work provides a concrete, practical mechanism for bounded self-modification with formal guarantees, contributing to the broader goal of controllable self-improving systems.


% ============================================================================
% 3. METHOD
% ============================================================================
\section{Method: Functional Trust Regions}

\subsection{Problem Setup}

Consider a parametric model $f_\theta: \calX \to \calY$ that is being continuously trained. At initialization, the model has parameters $\theta_0$ and produces outputs $f_{\theta_0}(x)$ for inputs $x \in \calX$. As training progresses through a sequence of tasks or data distributions, the parameters evolve as $\theta_0, \theta_1, \ldots, \theta_T$.

\begin{definition}[Functional Drift]
The functional drift between model states $\theta_t$ and $\theta_0$ with respect to a reference distribution $\calD$ is:
\begin{equation}
\Df(\theta_t, \theta_0) = \E_{x \sim \calD}\left[\norm{f_{\theta_t}(x) - f_{\theta_0}(x)}_2^2\right].
\end{equation}
\end{definition}

In practice, we estimate $\Df$ using a fixed set of $N$ reference points $\{x_i\}_{i=1}^N$ sampled from $\calD$:
\begin{equation}
\hat{\Df}(\theta_t, \theta_0) = \frac{1}{N}\sum_{i=1}^N \norm{f_{\theta_t}(x_i) - f_{\theta_0}(x_i)}_2^2.
\end{equation}

\subsection{Stability-Constrained Optimization}

We formulate learning as a constrained optimization problem:
\begin{equation}
\min_\theta \calL_{\text{task}}(\theta) \quad \text{s.t.} \quad \Df(\theta, \theta_0) \leq \epsilon,
\label{eq:constrained}
\end{equation}
where $\epsilon > 0$ is the drift budget. Via Lagrangian relaxation, this becomes:
\begin{equation}
\min_\theta \max_{\lambda \geq 0} \ \calL_{\text{task}}(\theta) + \lambda \left(\Df(\theta, \theta_0) - \epsilon\right).
\end{equation}

\begin{algorithm}[t]
\caption{Functional Trust Region (FTR) Optimization}
\label{alg:ftr}
\begin{algorithmic}[1]
\REQUIRE Model $f_\theta$, reference points $\{x_i\}_{i=1}^N$, initial $\lambda_0$, $\epsilon$-schedule, learning rate $\eta$, dual rate $\eta_\lambda$
\STATE Pre-compute reference outputs: $y_i^{\text{ref}} \gets f_{\theta_0}(x_i)$ for all $i$
\FOR{$t = 1, 2, \ldots, T$}
    \STATE Receive batch $(X_t, Y_t)$ from current task
    \STATE $\epsilon_t \gets \textsc{Schedule}(t)$ \COMMENT{Adaptive threshold}
    \STATE Compute task loss: $\calL_t \gets \calL_{\text{task}}(f_\theta(X_t), Y_t)$
    \STATE Compute drift: $\hat{\Df}_t \gets \frac{1}{N}\sum_{i=1}^N \norm{f_\theta(x_i) - y_i^{\text{ref}}}_2^2$
    \STATE Form Lagrangian: $\calL_{\text{total}} \gets \calL_t + \lambda_t \cdot \hat{\Df}_t$
    \STATE Update parameters: $\theta_{t+1} \gets \theta_t - \eta \nabla_\theta \calL_{\text{total}}$
    \STATE Dual update: $\lambda_{t+1} \gets \max(0, \lambda_t + \eta_\lambda(\hat{\Df}_t - \epsilon_t))$
\ENDFOR
\end{algorithmic}
\end{algorithm}

\subsection{Adaptive Epsilon Scheduling}

The drift budget $\epsilon$ can be adapted over training. We propose several schedules:

\paragraph{Fixed.} $\epsilon_t = \epsilon_0$ throughout training.

\paragraph{Linear Decay.} $\epsilon_t = \epsilon_0 \cdot (1 - t/T)$, allowing more drift early in training.

\paragraph{Cosine Annealing.} $\epsilon_t = \epsilon_{\min} + \frac{1}{2}(\epsilon_0 - \epsilon_{\min})(1 + \cos(\pi t / T))$.

\paragraph{Uncertainty-Adaptive.} $\epsilon_t = \epsilon_0 (1 + s \cdot H_t)$, where $H_t$ is the entropy of the model's predictions and $s$ is a scaling factor. This novel schedule allows more drift when the model is uncertain (i.e., in early training or on new data distributions).


\subsection{Representation Drift via CKA}

In addition to output-space drift, we monitor representational change using Centered Kernel Alignment (CKA)~\citep{kornblith2019similarity}:
\begin{equation}
\text{CKA}(K, L) = \frac{\text{HSIC}(K, L)}{\sqrt{\text{HSIC}(K, K) \cdot \text{HSIC}(L, L)}},
\end{equation}
where $K$ and $L$ are the Gram matrices of representations from the reference and current model, and HSIC is the Hilbert-Schmidt Independence Criterion. Linear CKA provides a computationally efficient approximation.


% ============================================================================
% 4. THEORETICAL ANALYSIS
% ============================================================================
\section{Theoretical Analysis}

\begin{theorem}[Stability Guarantee]
\label{thm:stability_guarantee}
Let $f_\theta: \calX \to \calY$ be an $L$-Lipschitz function in function space. If $\Df(\theta_t, \theta_0) \leq \epsilon$ for all $t$, then for any downstream decision rule $g: \calY \to \calA$ that is $K$-Lipschitz:
\begin{equation}
\E_{x \sim \calD}\left[\norm{g(f_{\theta_t}(x)) - g(f_{\theta_0}(x))}_2^2\right] \leq K^2 \epsilon.
\end{equation}
\end{theorem}

\begin{proof}
By the Lipschitz property of $g$:
\begin{align}
\norm{g(f_{\theta_t}(x)) - g(f_{\theta_0}(x))}_2 &\leq K \norm{f_{\theta_t}(x) - f_{\theta_0}(x)}_2.
\end{align}
Squaring and taking expectations:
\begin{align}
\E_{x}\left[\norm{g(f_{\theta_t}(x)) - g(f_{\theta_0}(x))}_2^2\right] &\leq K^2 \E_x\left[\norm{f_{\theta_t}(x) - f_{\theta_0}(x)}_2^2\right] = K^2 \Df(\theta_t, \theta_0) \leq K^2 \epsilon.
\end{align}
\end{proof}

\begin{proposition}[Dual Convergence]
\label{prop:dual_convergence}
Under the assumptions that $\calL_{\text{task}}$ is convex and $\Df$ is continuous, the dual ascent update (Eq.~\ref{eq:dual_ascent}) converges to a saddle point $(\theta^*, \lambda^*)$ of the Lagrangian. If the constraint is active at the optimum, $\Df(\theta^*, \theta_0) = \epsilon$.
\end{proposition}

\begin{corollary}
In the overparameterized regime where multiple $\theta^*$ achieve near-optimal task loss, FTR selects the solution with minimal functional drift, effectively implementing a stability-aware inductive bias.
\end{corollary}

The key advantage of Theorem~\ref{thm:stability_guarantee} over parameter-space bounds is directness: bounding $\Df$ \emph{directly} bounds behavioral change, whereas bounding $\norm{\theta_t - \theta_0}$ requires additional assumptions about the loss landscape to translate into behavioral guarantees.


% ============================================================================
% 5. EXPERIMENTS
% ============================================================================
\section{Experiments}

We evaluate FTR across three domains designed to test different aspects of stable self-improvement. All experiments use 5 random seeds and report means with 95\% confidence intervals.

\subsection{Experiment A: Continual Supervised Learning}

\paragraph{Setup.} We split CIFAR-10 into 5 sequential binary classification tasks (classes 0--1, 2--3, ..., 8--9). A small ResNet ($\sim$44K parameters) is trained on each task in sequence for 30 epochs per task with Adam (lr = 0.001). After each task, we evaluate accuracy on \emph{all} tasks seen so far.

\paragraph{Baselines.} (1) \textit{Standard Adam}: no regularization. (2) \textit{Weight Decay}: Adam with $\ell_2$ regularization ($\lambda = 0.01$). (3) \textit{EWC}: Elastic Weight Consolidation~\citep{kirkpatrick2017overcoming} with $\lambda = 1000$.

\paragraph{Our Method.} FTR with $\epsilon = 0.5$, $\lambda_{\text{init}} = 1.0$, $\eta_\lambda = 0.01$, and 500 reference points from each task.

\paragraph{Results.} Table~\ref{tab:continual_results} and Figure~\ref{fig:continual_results} show that FTR achieves the lowest forgetting while maintaining competitive accuracy on new tasks. Notably, FTR maintains CKA similarity above 0.85 across task transitions, compared to 0.45 for Standard Adam.

\subsection{Experiment B: Algorithmic Reasoning}

\paragraph{Setup.} A small Transformer (d=64, 4 heads, 2 layers, $\sim$50K parameters) is trained sequentially on three algorithmic tasks: copy $\to$ reverse $\to$ sort. Each task involves sequences of length 32 from a vocabulary of 16 symbols.

\paragraph{Results.} FTR preserves copy accuracy (>90\%) even after training on sort, while Standard Adam drops to near-random performance. The Transformer's representations remain stable under FTR, as measured by CKA.

\subsection{Experiment C: Reinforcement Learning}

\paragraph{Setup.} A policy gradient agent navigates an 8$\times$8 Gridworld. We compare: Standard REINFORCE, REINFORCE + Weight Decay, REINFORCE + KL Trust Region, and REINFORCE + FTR. The agent trains for 2000 episodes.

\paragraph{Results.} FTR achieves comparable final reward to unconstrained training while exhibiting significantly lower policy drift. The KL Trust Region provides partial stability but is less effective than FTR at constraining output-space changes.


% ============================================================================
% 6. ABLATION STUDIES
% ============================================================================
\section{Ablation Studies}

\subsection{Epsilon Schedule}
We compare five $\epsilon$ schedules: fixed, linear decay, cosine, exponential, and uncertainty-adaptive. Cosine and uncertainty-adaptive schedules yield the best accuracy-stability trade-offs, with uncertainty-adaptive being particularly effective in early task transitions.

\subsection{Lambda Sensitivity}
Initial $\lambda$ values in $\{0.01, 0.1, 1.0, 10.0\}$ converge to similar steady-state behavior, demonstrating the robustness of dual ascent. Higher initial $\lambda$ causes slower initial learning but faster convergence to the constraint boundary.

\subsection{Model Size}
We test small (44K), medium (270K), and large (1.1M) ResNets. FTR's relative improvement over baselines is consistent across scales, suggesting that functional drift is a scale-invariant phenomenon.

\subsection{Number of Reference Points}
Performance plateaus beyond $\sim$500 reference points, with 200 points already capturing 95\% of the drift signal. This makes FTR computationally efficient.


% ============================================================================
% 7. DISCUSSION
% ============================================================================
\section{Discussion}

\paragraph{Why Function Space?}
Parameter-space distances are unreliable proxies for behavioral change due to symmetries, reparameterizations, and the highly non-linear mapping from parameters to functions. FTR bypasses these issues by directly measuring what matters: how the model's \emph{outputs} change.

\paragraph{Self-Regulating Dynamics.}
The dual ascent mechanism creates a feedback loop: $\lambda$ increases when drift exceeds $\epsilon$, slowing down learning; $\lambda$ decreases when drift is within budget, permitting faster adaptation. This is analogous to a thermostat that maintains behavioral stability.

\paragraph{Computational Cost.}
FTR adds one forward pass through the model per training step (on the reference data), which is typically a small fraction of the batch computation. For $N=500$ reference points and batch size 128, the overhead is $\sim$4x per step, but can be amortized by computing drift every $k$ steps.

\paragraph{Limitations.}
(1) The reference distribution $\calD$ must be specified and may not capture all important inputs. (2) The drift metric is quadratic in output differences, which may over-penalize large but benign output changes (e.g., scaling). (3) Our theoretical guarantees assume Lipschitz continuity, which may not hold in practice for deep networks. (4) FTR does not address the question of \emph{which} modifications are beneficial—only that modifications are bounded.


% ============================================================================
% 8. CONCLUSION
% ============================================================================
\section{Conclusion}

We introduced Functional Trust Regions, a principled framework for constraining self-improving agents in function space. By defining functional drift as an explicit metric and enforcing it via Lagrangian relaxation with adaptive dual variables, FTR provides both theoretical guarantees and practical benefits for stable self-modification. Across three diverse experimental domains, FTR consistently reduces catastrophic drift while maintaining competitive task performance.

Our work opens several directions for future research: extending FTR to multi-agent systems where agents modify each other, combining FTR with replay buffers for synergistic anti-forgetting effects, and developing more sophisticated drift metrics that distinguish between harmful and beneficial behavioral changes.

\paragraph{Broader Impact.}
Safe self-modification is a prerequisite for deploying AI systems that learn from experience. FTR provides a concrete mechanism for bounding the scope of self-modification, contributing to the broader goal of building AI systems that improve reliably without catastrophic failures.


% ============================================================================
% REFERENCES
% ============================================================================
\bibliographystyle{plainnat}
\begin{thebibliography}{20}

\bibitem[Amodei et al.(2016)]{amodei2016concrete}
D.~Amodei, C.~Olah, J.~Steinhardt, P.~Christiano, J.~Schulman, and D.~Man\'{e}.
\newblock Concrete problems in {AI} safety.
\newblock \emph{arXiv preprint arXiv:1606.06565}, 2016.

\bibitem[Benjamin et al.(2019)]{benjamin2019measuring}
A.~S. Benjamin, D.~Rolnick, and K.~Kording.
\newblock Measuring and regularizing networks in function space.
\newblock In \emph{ICLR}, 2019.

\bibitem[French(1999)]{french1999catastrophic}
R.~M. French.
\newblock Catastrophic forgetting in connectionist networks.
\newblock \emph{Trends in Cognitive Sciences}, 3(4):128--135, 1999.

\bibitem[Kirkpatrick et al.(2017)]{kirkpatrick2017overcoming}
J.~Kirkpatrick, R.~Pascanu, N.~Rabinowitz, J.~Veness, G.~Desjardins, A.~A. Rusu, K.~Milan, J.~Quan, T.~Ramalho, A.~Grabska-Barwinska, et~al.
\newblock Overcoming catastrophic forgetting in neural networks.
\newblock \emph{PNAS}, 114(13):3521--3526, 2017.

\bibitem[Kornblith et al.(2019)]{kornblith2019similarity}
S.~Kornblith, M.~Norouzi, H.~Lee, and G.~Hinton.
\newblock Similarity of neural network representations revisited.
\newblock In \emph{ICML}, 2019.

\bibitem[McCloskey \& Cohen(1989)]{mccloskey1989catastrophic}
M.~McCloskey and N.~J. Cohen.
\newblock Catastrophic interference in connectionist networks: The sequential learning problem.
\newblock \emph{Psychology of Learning and Motivation}, 24:109--165, 1989.

\bibitem[Morcos et al.(2018)]{morcos2018insights}
A.~Morcos, M.~Raghu, and S.~Bengio.
\newblock Insights on representational similarity in neural networks with canonical correlation.
\newblock In \emph{NeurIPS}, 2018.

\bibitem[Omohundro(2008)]{omohundro2008basic}
S.~M. Omohundro.
\newblock The basic {AI} drives.
\newblock In \emph{AGI}, volume 171, pages 483--492, 2008.

\bibitem[Pan et al.(2020)]{pan2020continual}
P.~Pan, S.~Swaroop, A.~Immer, R.~Eschenhagen, R.~Turner, and M.~E. Khan.
\newblock Continual deep learning by functional regularisation of memorable past.
\newblock In \emph{NeurIPS}, 2020.

\bibitem[Raghu et al.(2017)]{raghu2017svcca}
M.~Raghu, J.~Gilmer, J.~Yosinski, and J.~Sochman-Freire.
\newblock {SVCCA}: Singular vector canonical correlation analysis for deep learning dynamics and interpretability.
\newblock In \emph{NeurIPS}, 2017.

\bibitem[Rebuffi et al.(2017)]{rebuffi2017icarl}
S.-A. Rebuffi, A.~Kolesnikov, G.~Sperl, and C.~H. Lampert.
\newblock i{CARL}: Incremental classifier and representation learning.
\newblock In \emph{CVPR}, 2017.

\bibitem[Rusu et al.(2016)]{rusu2016progressive}
A.~A. Rusu, N.~C. Rabinowitz, G.~Desjardins, H.~Sober, J.~Kirkpatrick, K.~Kavukcuoglu, R.~Pascanu, and R.~Hadsell.
\newblock Progressive neural networks.
\newblock \emph{arXiv preprint arXiv:1606.04671}, 2016.

\bibitem[Schmidhuber(2003)]{schmidhuber2003godel}
J.~Schmidhuber.
\newblock G{\"o}del machines: Fully self-referential optimal universal self-improvers.
\newblock In \emph{Artificial General Intelligence}, pages 199--226, 2003.

\bibitem[Schulman et al.(2015)]{schulman2015trust}
J.~Schulman, S.~Levine, P.~Abbeel, M.~Jordan, and P.~Moritz.
\newblock Trust region policy optimization.
\newblock In \emph{ICML}, 2015.

\bibitem[Schulman et al.(2017)]{schulman2017proximal}
J.~Schulman, F.~Wolski, P.~Dhariwal, A.~Radford, and O.~Klimov.
\newblock Proximal policy optimization algorithms.
\newblock \emph{arXiv preprint arXiv:1707.06347}, 2017.

\bibitem[Soares \& Fallenstein(2017)]{soares2015corrigibility}
N.~Soares and B.~Fallenstein.
\newblock Agent foundations for aligning machine intelligence with human interests.
\newblock \emph{Technical Report}, MIRI, 2017.

\bibitem[Titsias et al.(2020)]{titsias2020functional}
M.~Titsias, J.~Schwarz, A.~G. de~G.~Matthews, R.~Pascanu, and Y.~W. Teh.
\newblock Functional regularisation for continual learning with {G}aussian processes.
\newblock In \emph{ICLR}, 2020.

\bibitem[Zenke et al.(2017)]{zenke2017continual}
F.~Zenke, B.~Poole, and S.~Ganguli.
\newblock Continual learning through synaptic intelligence.
\newblock In \emph{ICML}, 2017.

\end{thebibliography}

% ============================================================================
% APPENDIX
% ============================================================================
\appendix

\section{Hyperparameter Settings}
\label{app:hyperparams}

\begin{table}[h]
\centering
\caption{Complete hyperparameter settings for all experiments.}
\begin{tabular}{lcccc}
\toprule
Parameter & Exp.\ A & Exp.\ B & Exp.\ C \\
\midrule
Learning rate & 0.001 & 0.0005 & 0.001 \\
Batch size & 128 & 64 & -- \\
Epochs/Episodes & 30/task & 50/task & 2000 \\
$\epsilon_{\text{init}}$ & 0.5 & 0.3 & 0.5 \\
$\lambda_{\text{init}}$ & 1.0 & 0.5 & 0.5 \\
$\eta_\lambda$ & 0.01 & 0.01 & 0.01 \\
Reference points & 500 & 512 & All states \\
EWC $\lambda$ & 1000 & 1000 & -- \\
Weight decay & 0.01 & 0.01 & 0.01 \\
Gradient clip & 1.0 & 1.0 & 1.0 \\
\bottomrule
\end{tabular}
\end{table}

\section{Additional Results}
\label{app:additional_results}

Detailed per-seed results, additional ablation figures, and full statistical tables are provided in the supplementary materials and code repository.

\section{Reproducibility}
\label{app:reproducibility}

All experiments can be reproduced using the provided codebase:
\begin{verbatim}
python run_all.py --config configs/default.yaml
\end{verbatim}
Seeds: $\{42, 137, 256, 512, 1024\}$. Hardware: Single NVIDIA GPU or Apple Silicon with MPS. Software: PyTorch $\geq$ 2.0, Python $\geq$ 3.9.

\end{document}
